\documentclass[11pt,a4paper]{article}

\usepackage[margin=2.4cm]{geometry}
\usepackage{amsmath,amssymb,amsthm}
\usepackage{booktabs}
\usepackage{array}
\usepackage{longtable}
\usepackage{listings}
\usepackage{xcolor}
\usepackage{enumitem}

\definecolor{BlueLink}{RGB}{20,80,160}
\definecolor{CodeBg}{RGB}{247,247,250}
\definecolor{CodeKw}{RGB}{140,20,120}
\definecolor{CodeStr}{RGB}{20,110,60}
\definecolor{CodeCom}{RGB}{110,110,120}

\usepackage[colorlinks=true,linkcolor=BlueLink,citecolor=BlueLink,urlcolor=BlueLink]{hyperref}

\lstdefinestyle{py}{
  language=Python,
  backgroundcolor=\color{CodeBg},
  basicstyle=\ttfamily\footnotesize,
  keywordstyle=\color{CodeKw}\bfseries,
  stringstyle=\color{CodeStr},
  commentstyle=\color{CodeCom}\itshape,
  numbers=left, numberstyle=\tiny\color{CodeCom}, numbersep=7pt,
  frame=single, framesep=4pt, rulecolor=\color{CodeBg!50!gray},
  showstringspaces=false, breaklines=true, columns=fullflexible,
  captionpos=b, xleftmargin=14pt,
}
\lstset{style=py}

% a little extra give, so the unbreakable \shape{...} and \code{...} tokens
% cannot push a line past the margin
\emergencystretch=2.5em

% shape macro: \shape{B,N,3}
\newcommand{\shape}[1]{\ensuremath{[\,\texttt{#1}\,]}}
\newcommand{\code}[1]{\texttt{#1}}
\newcommand{\R}{\mathbb{R}}
\newcommand{\E}{\mathbb{E}}
\newcommand{\bx}{\mathbf{x}}
\newcommand{\bh}{\mathbf{h}}
\newcommand{\bz}{\mathbf{z}}
\newcommand{\bm}{\boldsymbol{\epsilon}}
\newcommand{\file}[1]{\textcolor{BlueLink}{\texttt{#1}}}

\title{\textbf{Equivariant Diffusion for 3D Molecule Generation on QM9}\\[4pt]
\large A line-by-line walkthrough of the implementation}
\author{Code documentation}
\date{}

\begin{document}
\maketitle

\begin{abstract}
\noindent
This document explains an implementation of an $E(3)$-equivariant denoising
diffusion probabilistic model that generates small organic molecules as 3D point
clouds --- atom types \emph{and} Cartesian coordinates jointly --- trained on
QM9. It is a compact reimplementation of EDM (Hoogeboom et al., 2022) with an
EGNN (Satorras et al., 2021) denoiser, depending only on \code{torch} and
\code{numpy}.

The emphasis throughout is on \textbf{tensor shapes}: every array is annotated
with its dimensions, and each transformation states what the shape was before
and after. All shapes, parameter counts and schedule values quoted here were
measured by instrumenting the running code, not derived on paper.
\end{abstract}

\tableofcontents

\section{Notation and the dimension glossary}

Six symbols determine every shape in the codebase. Fix them once and the rest
of the document follows.

\begin{center}
\begin{tabular}{@{}clll@{}}
\toprule
\textbf{Symbol} & \textbf{Meaning} & \textbf{Code name} & \textbf{Typical value} \\
\midrule
$B$ & batch size (molecules per batch)   & \code{batch\_size} & $64$ \\
$N$ & padded atom count (max in split)   & \code{max\_n}      & $29$ with H, $9$ without \\
$K$ & number of atom types               & \code{num\_types}  & $5$ (H,C,N,O,F); $4$ without H \\
$H$ & hidden width of the EGNN           & \code{hidden}      & $192$ \\
$L$ & number of EGNN layers              & \code{n\_layers}   & $6$ \\
$T$ & diffusion timesteps                & \code{timesteps}   & $1000$ \\
\bottomrule
\end{tabular}
\end{center}

\noindent
A molecule is represented \emph{densely}: every molecule in a batch is padded to
$N$ atoms and carries a binary mask marking which slots are real.

\begin{center}
\begin{tabular}{@{}lll@{}}
\toprule
\textbf{Tensor} & \textbf{Shape} & \textbf{Contents} \\
\midrule
\code{x}    & \shape{B,N,3} & atom coordinates in \r{A}ngstr\"om, zeros in padded slots \\
\code{h}    & \shape{B,N,K} & one-hot atom types, all-zero rows in padded slots \\
\code{mask} & \shape{B,N}   & $1.0$ for a real atom, $0.0$ for padding \\
\bottomrule
\end{tabular}
\end{center}

\noindent
The mask is the load-bearing part of the dense design. It appears in every
sum, every mean and every layer output; a single missing multiplication by it
lets padded slots contribute to a molecule's chemistry. Section~\ref{sec:masking}
collects the places where this matters.

\paragraph{Why $N$ differs between modes.} $N$ is not a constant of the model
but of the \emph{data split}: \code{QM9Dataset} truncates its padded arrays to
the largest molecule it actually contains. With hydrogens that is $29$; with
\code{--remove-h} the heavy-atom skeletons top out around $9$. Since the cost of
a forward pass grows as $N^2$ (Section~\ref{sec:cost}), this single number
accounts for the $\sim\!6\times$ speed difference between the two modes.

\section{The data pipeline: \file{qm9\_diffusion/data.py}}

\subsection{From the raw release to a padded cache}

QM9 is 133{,}885 small organic molecules over $\{$H, C, N, O, F$\}$, each with a
DFT-relaxed geometry. The loader needs only types and coordinates, so it
discards everything else and caches one \code{.npz}:

\begin{center}
\begin{tabular}{@{}llll@{}}
\toprule
\textbf{Stage} & \textbf{Representation} & \textbf{Shape} & \textbf{dtype} \\
\midrule
raw            & \code{gdb9.sdf}, one text record per molecule & --- & text \\
parsed         & Python list of per-molecule arrays  & $M\times(n_i,3)$ & \code{float32} \\
cached         & \code{positions}   & \shape{133885,29,3} & \code{float32} \\
               & \code{atom\_types} & \shape{133885,29}   & \code{int8} \\
               & \code{num\_atoms}  & \shape{133885}      & \code{int16} \\
\bottomrule
\end{tabular}
\end{center}

\noindent
Molecule $i$ occupies rows $0\ldots n_i-1$ of its slice; rows $n_i\ldots 28$ are
zero. The cache is $\approx 50$\,MB and is built once.

\paragraph{The V2000 atom line.} \code{\_parse\_sdf} reads coordinates by fixed
column offsets rather than by splitting on whitespace, because the MOL format is
column-oriented and the fields can abut:
\begin{lstlisting}[caption={Fixed-width parse of one SDF atom line \code{data.py}}]
sym = line[31:34].strip()                       # element symbol
pos[i] = (float(line[0:10]),                    # x, columns  1-10
          float(line[10:20]),                   # y, columns 11-20
          float(line[20:30]))                    # z, columns 21-30
\end{lstlisting}
That the resulting geometries reproduce the published stability statistics
exactly (Section~\ref{sec:eval}) is the evidence that this parse is correct.

\subsection{What one training example looks like}

\code{QM9Dataset.\_\_getitem\_\_} converts stored integer types into the dense
masked form the model consumes:

\begin{lstlisting}[caption={\code{QM9Dataset.\_\_getitem\_\_} --- integer types to masked one-hot}]
n = int(self.num_atoms[i])          # scalar: true atom count
mask = torch.zeros(self.max_n)      # [N]
mask[:n] = 1.0                      # [N],   n ones then zeros
one_hot = F.one_hot(self.atom_types[i], self.num_types).float()   # [N] -> [N,K]
return {
    "x":    self.positions[i] * mask[:, None],   # [N,3] * [N,1] -> [N,3]
    "h":    one_hot           * mask[:, None],   # [N,K] * [N,1] -> [N,K]
    "mask": mask,                                # [N]
}
\end{lstlisting}

\noindent
Two shape details deserve attention:
\begin{itemize}[leftmargin=1.4em]
  \item \code{mask[:, None]} promotes \shape{N} to \shape{N,1} so it broadcasts
        against both \shape{N,3} and \shape{N,K}. The same idiom recurs
        throughout the model as \code{node\_mask.unsqueeze(-1)}.
  \item The one-hot of a padded slot would be $e_0$ (type index $0$), i.e.\ a
        spurious hydrogen. Multiplying by the mask replaces it with the all-zero
        row, which is what makes padded atoms genuinely typeless.
\end{itemize}

\noindent
The \code{DataLoader} then stacks $B$ such dictionaries, prepending the batch
axis and producing exactly the \shape{B,N,3}, \shape{B,N,K}, \shape{B,N} triple
of Section~1.

\subsection{The size distribution $p(N_{\text{atoms}})$}

At generation time the model must be told \emph{how many} atoms to generate --- the
diffusion process itself does not create or destroy nodes. The empirical
histogram supplies that:

\begin{lstlisting}[caption={\code{size\_histogram} --- an empirical prior over molecule size}]
hist = torch.bincount(self.num_atoms, minlength=self.max_n + 1).float()  # [N+1]
return hist / hist.sum()                                                 # [N+1]
\end{lstlisting}

\noindent
Index $k$ holds $p(n_{\text{atoms}}=k)$, hence length $N+1$ rather than $N$.
Sampling proceeds by drawing sizes and turning them into a mask:
\begin{lstlisting}
sizes = torch.multinomial(size_hist, n_samples, replacement=True)   # [B]
mask  = (torch.arange(max_n)[None, :] < sizes[:, None]).float()     # [1,N] vs [B,1] -> [B,N]
\end{lstlisting}
The comparison broadcasts a row vector of positions against a column vector of
lengths, yielding a left-aligned boolean mask in one operation.

\section{The zero-centre-of-mass subspace}
\label{sec:com}

This is the one piece of geometry that must be understood before the model makes
sense, so it comes before the network.

\paragraph{The problem.} We want the learned density $p(\bx)$ over
$\bx\in\R^{n\times 3}$ to be invariant to rigid motions: a molecule and its
translated copy are the same molecule. But no probability density on $\R^{n\times 3}$
can be translation invariant. If $p(\bx)=p(\bx+t\mathbf{1})$ for every
$t\in\R^3$, then $p$ is constant along $3$ unbounded directions and cannot
integrate to $1$.

\paragraph{The fix.} Quotient out the offending directions. Work on the linear
subspace of configurations whose centre of mass is the origin,
\begin{equation}
  \mathcal{X}_0 \;=\; \Bigl\{\, \bx \in \R^{n\times3} \;:\; \textstyle\sum_{i=1}^{n} \bx_i = \mathbf{0} \,\Bigr\},
  \qquad \dim \mathcal{X}_0 \;=\; 3(n-1).
\end{equation}
A density on $\mathcal{X}_0$ \emph{is} normalisable, and pushing it forward by
"place the centre of mass anywhere" gives the translation-invariant object we
actually want. Every coordinate quantity in the code therefore lives in
$\mathcal{X}_0$: the data, the noise, and the network's coordinate output.

\paragraph{The projection.} It is a mean subtraction, but a \emph{masked} one:
the mean must be taken over the $n_b$ real atoms of molecule $b$, not over all
$N$ padded slots.

\begin{lstlisting}[caption={\code{remove\_mean} --- projection onto $\mathcal{X}_0$ \code{diffusion.py}}]
def remove_mean(x, node_mask):                       # x [B,N,3], node_mask [B,N]
    nm   = node_mask.unsqueeze(-1)                   # [B,N]   -> [B,N,1]
    n    = node_mask.sum(1).view(-1, 1, 1).clamp(min=1.0)   # [B,N] -> [B] -> [B,1,1]
    mean = (x * nm).sum(1, keepdim=True) / n          # [B,N,3] -> [B,1,3]
    return (x - mean) * nm                            # [B,1,3] broadcasts -> [B,N,3]
\end{lstlisting}

\noindent
Shape trace: \shape{B,N,3} $\to$ sum over the atom axis $\to$ \shape{B,1,3}
$\to$ broadcast back to \shape{B,N,3}. The \code{keepdim=True} is what makes the
subtraction broadcast without an explicit \code{unsqueeze}. The
\code{clamp(min=1.0)} guards against a division by zero for an all-padding row.

\paragraph{Noise must be projected too.} Gaussian noise added to a point in
$\mathcal{X}_0$ has to stay in $\mathcal{X}_0$, so the sampler never uses raw
\code{randn} for coordinates:
\begin{lstlisting}
def sample_zero_com_noise(shape, node_mask, device):
    return remove_mean(torch.randn(shape, device=device), node_mask)   # [B,N,3]
\end{lstlisting}
This projected noise is \emph{not} standard normal: it has variance
$\tfrac{n-1}{n}$ per coordinate, a fact used to check the sampler in
Section~\ref{sec:calib}.

\section{The denoiser: \file{qm9\_diffusion/egnn.py}}
\label{sec:egnn}

The network's job is a regression: given a noisy molecule and a timestep,
predict the noise that was added. It returns two tensors,
$\hat{\bm}_x$ \shape{B,N,3} and $\hat{\bm}_h$ \shape{B,N,K}.

\subsection{Why an EGNN and not an MLP}

Rotating a molecule rotates the noise that was added to it. A generic network
would have to learn this from data, spending capacity on a symmetry we already
know. An EGNN builds it in. Writing $R$ for a rotation/reflection in $O(3)$, the
network satisfies
\begin{equation}
  \hat{\bm}_x(\bx R,\,\bh,\,t) \;=\; \hat{\bm}_x(\bx,\,\bh,\,t)\,R,
  \qquad
  \hat{\bm}_h(\bx R,\,\bh,\,t) \;=\; \hat{\bm}_h(\bx,\,\bh,\,t),
\end{equation}
i.e.\ the coordinate output is \emph{equivariant} and the type output is
\emph{invariant}. The two structural facts that produce this:
\begin{enumerate}[leftmargin=1.6em]
  \item messages are functions of squared interatomic distances
        $\lVert \bx_i-\bx_j\rVert^2$, which are $O(3)$-invariant;
  \item coordinates are updated only along relative-position vectors
        $\bx_i-\bx_j$, which are $O(3)$-equivariant.
\end{enumerate}
Scalars times equivariant vectors stay equivariant, and the composition of
equivariant layers is equivariant, so the property holds for the whole stack.
\code{test\_smoke.py} verifies it numerically to $5\times10^{-7}$.

\subsection{Input embedding: where the timestep enters}

\begin{lstlisting}[caption={\code{EGNNDynamics.forward} --- masks and embedding}]
b, n, _ = x.shape                                  # B, N
nm  = node_mask.unsqueeze(-1)                      # [B,N] -> [B,N,1]
eye = torch.eye(n).view(1, n, n, 1)                # [N,N] -> [1,N,N,1]
edge_mask = (nm.unsqueeze(2) * nm.unsqueeze(1)) * (1.0 - eye)   # -> [B,N,N,1]

t_node = t.view(b, 1, 1).expand(-1, n, 1)          # [B] -> [B,1,1] -> [B,N,1]
hidden = self.embed(torch.cat([h, t_node], -1)) * nm   # [B,N,K+1] -> [B,N,H]
\end{lstlisting}

\noindent
Three shape manipulations, each worth naming:

\begin{center}\small
\begin{tabular}{@{}lccp{5.1cm}@{}}
\toprule
\textbf{Expression} & \textbf{Before} & \textbf{After} & \textbf{Purpose} \\
\midrule
\code{nm.unsqueeze(2)}\newline\code{\ \ \ \ * nm.unsqueeze(1)} & \shape{B,N,1} & \shape{B,N,N,1} & outer product of the masks: edge $(i,j)$ is live iff both ends are real \\
\addlinespace
\code{* (1.0 - eye)} & \shape{B,N,N,1} & \shape{B,N,N,1} & delete the self-loops \\
\addlinespace
\code{t.view(b,1,1)}\newline\code{\ \ \ \ .expand(-1,n,1)} & \shape{B} & \shape{B,N,1} & broadcast the scalar time to every node \\
\bottomrule
\end{tabular}
\end{center}

\noindent
The time conditioning is deliberately minimal: $t$ is a single scalar in $[0,1]$
concatenated as one extra input channel, so the embedding is a linear map
$\R^{K+1}\!\to\!\R^{H}$ ($6\to192$ by default, only $1{,}344$ parameters). There
is no sinusoidal Fourier embedding here --- a design simplification relative to
image diffusion models, where the timestep is usually embedded richly.

\code{expand} rather than \code{repeat} is used because it creates a
\emph{view} with stride $0$ along the atom axis; no $B\times N$ block of memory
is written for what is logically one number per molecule.

\subsection{One EGNN layer, shape by shape}

This is the core of the model. The layer maps
$(\bh,\bx)\in$\shape{B,N,H}$\times$\shape{B,N,3} to a tensor pair of the same
two shapes, but the intermediate quantities live on \emph{edges} and carry a
fourth axis. With $B=64$, $N=29$, $H=192$ the trace is:

\begin{center}
\begin{longtable}{@{}llll@{}}
\toprule
\textbf{Line} & \textbf{Quantity} & \textbf{Shape} & \textbf{Concrete} \\
\midrule
\endfirsthead
\toprule
\textbf{Line} & \textbf{Quantity} & \textbf{Shape} & \textbf{Concrete} \\
\midrule
\endhead
in  & \code{h}                & \shape{B,N,H}    & $(64,29,192)$ \\
in  & \code{x}                & \shape{B,N,3}    & $(64,29,3)$ \\
\addlinespace
1   & \code{diff}             & \shape{B,N,N,3}  & $(64,29,29,3)$ \\
2   & \code{d2}               & \shape{B,N,N,1}  & $(64,29,29,1)$ \\
3   & \code{norm}             & \shape{B,N,N,1}  & $(64,29,29,1)$ \\
\addlinespace
4   & \code{hi} (expand)      & \shape{B,N,N,H}  & $(64,29,29,192)$ \\
5   & \code{hj} (expand)      & \shape{B,N,N,H}  & $(64,29,29,192)$ \\
6   & \code{edge\_mlp} input  & \shape{B,N,N,2H{+}1} & $(64,29,29,385)$ \\
7   & \code{m} (message)      & \shape{B,N,N,H}  & $(64,29,29,192)$ \\
8   & \code{att} (gate)       & \shape{B,N,N,1}  & $(64,29,29,1)$ \\
9   & \code{m*att*edge\_mask} & \shape{B,N,N,H}  & $(64,29,29,192)$ \\
\addlinespace
10  & \code{coord\_mlp(m)}    & \shape{B,N,N,1}  & $(64,29,29,1)$ \\
11  & \code{trans}            & \shape{B,N,N,3}  & $(64,29,29,3)$ \\
12  & \code{n\_edges}         & \shape{B,N,1}    & $(64,29,1)$ \\
13  & \code{agg}              & \shape{B,N,H}    & $(64,29,192)$ \\
14  & \code{node\_mlp} input  & \shape{B,N,2H}   & $(64,29,384)$ \\
\addlinespace
out & \code{h}                & \shape{B,N,H}    & $(64,29,192)$ \\
out & \code{x}                & \shape{B,N,3}    & $(64,29,3)$ \\
\bottomrule
\end{longtable}
\end{center}

\noindent
The pattern is a round trip: \textbf{node axis $\to$ edge axis $\to$ node axis}.
Lines 1--9 lift \shape{B,N,$\cdot$} to \shape{B,N,N,$\cdot$}; lines 11--13
contract the third axis away with a sum. The code:

\begin{lstlisting}[caption={\code{EquivariantConv.forward}}]
diff = x.unsqueeze(2) - x.unsqueeze(1)     # [B,N,1,3] - [B,1,N,3] -> [B,N,N,3]
d2   = (diff**2).sum(-1, keepdim=True)     #                       -> [B,N,N,1]
norm = (d2 + 1e-8).sqrt()                  #                       -> [B,N,N,1]

hi = h.unsqueeze(2).expand(-1, -1, n, -1)  # [B,N,1,H] -> [B,N,N,H]   (sender)
hj = h.unsqueeze(1).expand(-1, n, -1, -1)  # [B,1,N,H] -> [B,N,N,H]   (receiver)
m  = self.edge_mlp(torch.cat([hi, hj, d2], dim=-1))   # [B,N,N,2H+1] -> [B,N,N,H]
m  = m * self.att_mlp(m) * edge_mask       # [B,N,N,H] * [B,N,N,1] * [B,N,N,1]

trans   = diff / (norm + 1.0) * torch.tanh(self.coord_mlp(m)) * self.coords_range
trans   = trans * edge_mask                          # [B,N,N,3]
n_edges = edge_mask.sum(dim=2).clamp(min=1.0)        # [B,N,N,1] -> [B,N,1]
x = x + (trans.sum(dim=2) / n_edges) * node_mask     # [B,N,N,3] -> [B,N,3]

agg = m.sum(dim=2) / n_edges                         # [B,N,N,H] -> [B,N,H]
h   = h + self.node_mlp(torch.cat([h, agg], dim=-1)) # [B,N,2H] -> [B,N,H]
return h * node_mask, x * node_mask
\end{lstlisting}

\paragraph{The broadcast that builds the edge axis.} \code{x.unsqueeze(2)} is
\shape{B,N,1,3} and \code{x.unsqueeze(1)} is \shape{B,1,N,3}; subtracting them
broadcasts both to \shape{B,N,N,3}, and entry $(b,i,j)$ of the result holds
$\bx_i-\bx_j$. This one line materialises all $N^2$ ordered pairs. The same
asymmetry names the two expanded copies of \code{h}: \code{hi} varies along
axis~1 and \code{hj} along axis~2, so the concatenation at position $(i,j)$ is
$[\bh_i,\bh_j,d^2_{ij}]$.

\paragraph{The equations.} In standard EGNN notation, with $\phi_e,\phi_h,\phi_x$
the three MLPs and $\phi_a$ the gate:
\begin{align}
  \mathbf{m}_{ij} &= \phi_a\!\bigl(\tilde{\mathbf{m}}_{ij}\bigr)\,\tilde{\mathbf{m}}_{ij}\,,
    \qquad \tilde{\mathbf{m}}_{ij} = \phi_e\bigl(\bh_i,\bh_j,\lVert\bx_i-\bx_j\rVert^2\bigr), \\
  \bx_i' &= \bx_i + \frac{1}{|\mathcal{N}_i|}\sum_{j\neq i}
            \frac{\bx_i-\bx_j}{\lVert\bx_i-\bx_j\rVert+1}\;
            \gamma\tanh\!\bigl(\phi_x(\mathbf{m}_{ij})\bigr), \\
  \bh_i' &= \bh_i + \phi_h\Bigl(\bh_i,\;\tfrac{1}{|\mathcal{N}_i|}\textstyle\sum_{j\neq i}\mathbf{m}_{ij}\Bigr).
\end{align}

\paragraph{Design choices inside the layer.}
\begin{description}[leftmargin=0pt,style=nextline]
\item[The $+1$ in the denominator.] \code{diff / (norm + 1.0)} is not a unit
      vector. Dividing by $\lVert\bx_i-\bx_j\rVert$ alone would make the update
      blow up for nearby atoms; the offset bounds the factor while keeping
      direction, and it is still equivariant because it scales an equivariant
      vector by an invariant scalar.
\item[$\gamma\tanh(\cdot)$.] \code{coords\_range} is passed as
      $\gamma=15/L=2.5$\,\r{A}, so a single layer can displace an atom by at most
      $2.5$\,\r{A} and the whole stack by at most $15$\,\r{A}. Without this bound
      the geometry can diverge early in training, and because distances feed
      back into the next layer's messages, a diverged geometry produces garbage
      messages rather than merely a large loss.
\item[Small initialisation of the coordinate head.]
      \code{xavier\_uniform\_(..., gain=0.001)} makes each layer start as
      approximately the identity on coordinates, so the model begins near
      "predict no displacement" and learns geometry gradually.
\item[Mean rather than sum aggregation.] Dividing by \code{n\_edges} keeps
      activation magnitudes independent of molecule size. Because
      \code{n\_edges} is computed from \code{edge\_mask}, the denominator is the
      true neighbour count $n_b-1$, not $N-1$: a $5$-atom molecule in a batch
      padded to $29$ is normalised by $4$.
\item[The gate.] \code{att\_mlp} is a $\text{Linear}(H,1)+\sigma$ producing one
      scalar in $(0,1)$ per edge. All $N^2$ pairs are present in a dense graph,
      most of them non-bonded; the gate lets the network suppress them instead of
      being forced to attend to everything.
\end{description}

\paragraph{A pitfall that this code hit.} Line 3 reads
\code{(d2 + 1e-8).sqrt()} rather than \code{d2.sqrt()}. Self-pairs ($i=j$) and
padding-padding pairs both have \code{diff}$=0$, hence \code{d2}$=0$, and
$\frac{d}{du}\sqrt{u}\to\infty$ as $u\to0$. The forward pass looks fine --- those
entries are killed by \code{edge\_mask} --- but the backward pass computes
$\infty \times 0 = \text{NaN}$ and the loss becomes NaN on the first step.
Masking an output does not protect its gradient; the guard has to be inside the
square root.

\subsection{Output heads}

\begin{lstlisting}[caption={From the last layer to $(\hat{\bm}_x,\hat{\bm}_h)$}]
eps_x = (x - x0) * nm                                       # [B,N,3]
eps_x = eps_x - eps_x.sum(1, keepdim=True) / node_mask.sum(1).view(b,1,1)
eps_x = eps_x * nm                                          # [B,N,3], zero-CoM
eps_h = self.decode(hidden) * nm                            # [B,N,H] -> [B,N,K]
\end{lstlisting}

\noindent
Two things happen here.
\begin{enumerate}[leftmargin=1.6em]
\item \textbf{The coordinate prediction is a displacement}, $\bx^{(L)}-\bx^{(0)}$,
      not the final coordinates. This is what makes the output transform
      correctly: both terms rotate together, so their difference does, whereas an
      absolute position would also have to track the input's placement.
\item \textbf{It is re-projected onto $\mathcal{X}_0$.} The target noise lives in
      the subspace, so the prediction must too --- otherwise the model wastes
      capacity on a component that can never reduce the loss, and the sampler
      accumulates centre-of-mass drift over $T$ steps. Verified to
      $1.5\times10^{-11}$ in the tests.
\end{enumerate}

\noindent
Note $\hat{\bm}_h$ is a free real-valued tensor in $\R^{B\times N\times K}$, not
a probability vector: there is no softmax. It predicts Gaussian noise on the
relaxed one-hot representation, which is why the loss below is a plain MSE.

\subsection{Parameter budget}

Measured for $H=192$, $L=6$, $K=5$:

\begin{center}
\begin{tabular}{@{}llrr@{}}
\toprule
\textbf{Module} & \textbf{Shape of the map} & \textbf{Params} & \textbf{$\times L$} \\
\midrule
\code{embed}     & $K{+}1 \to H$ \; $(6\to192)$              & $1{,}344$   & --- \\
\code{edge\_mlp} & $2H{+}1 \to H \to H$ \; $(385\to192\to192)$ & $111{,}168$ & $667{,}008$ \\
\code{att\_mlp}  & $H \to 1$                                  & $193$       & $1{,}158$ \\
\code{node\_mlp} & $2H \to H \to H$ \; $(384\to192\to192)$     & $110{,}976$ & $665{,}856$ \\
\code{coord\_mlp}& $H \to H \to 1$ \; $(192\to192\to1)$        & $37{,}249$  & $223{,}494$ \\
\code{decode}    & $H \to H \to K$ \; $(192\to192\to5)$        & $38{,}021$  & --- \\
\midrule
\textbf{Total}   &                                            & \multicolumn{2}{r}{$\mathbf{1{,}596{,}881}$} \\
\bottomrule
\end{tabular}
\end{center}

\noindent
$97.5\%$ of the parameters sit in the $L$ repeated layers, split almost evenly
between \code{edge\_mlp} and \code{node\_mlp}. The diffusion wrapper adds no
parameters --- its schedules are registered as \emph{buffers}, so they move with
\code{.to(device)} and are saved in checkpoints but are not optimised.

\section{The diffusion model: \file{qm9\_diffusion/diffusion.py}}

\subsection{The forward (noising) process}

Write $\bz_0=[\bx,\bh]$ for a clean molecule. The forward process is the usual
variance-preserving Gaussian diffusion,
\begin{equation}
  q(\bz_t \mid \bz_0) \;=\; \mathcal{N}\bigl(\bz_t;\;\alpha_t \bz_0,\;\sigma_t^2 \mathbf{I}\bigr),
  \qquad\text{implemented as}\qquad
  \bz_t \;=\; \alpha_t \bz_0 + \sigma_t \bm,\quad \bm\sim\mathcal{N}(0,\mathbf{I}),
\end{equation}
with $\alpha_t^2+\sigma_t^2=1$, so $\alpha_t$ decays from $\approx1$ to
$\approx0$ as $t: 0\to T$. Coordinates and types are noised jointly with the
\emph{same} $(\alpha_t,\sigma_t)$, but the coordinate noise is drawn in
$\mathcal{X}_0$ and the type noise in the full $\R^{K}$.

\subsection{The noise schedule}

\code{polynomial\_schedule} returns the array $\bar\alpha_t=\alpha_t^2$ for
$t=0,\dots,T$ --- length $T+1$, i.e.\ $1001$ entries for $T=1000$:
\begin{equation}
  \bar\alpha_t \;=\; (1-2s)\cdot\text{clip}\Bigl[\bigl(1-(t/T')^{2}\bigr)^{2}\Bigr] + s,
  \qquad s=10^{-5},\; T'=T+1 .
\end{equation}
\code{\_clip\_noise\_schedule} bounds the per-step ratio
$\bar\alpha_t/\bar\alpha_{t-1}\ge 0.001$ so that no single step annihilates the
signal, then re-accumulates with \code{cumprod}. The affine map by $s$ keeps
$\bar\alpha$ strictly inside $(0,1)$: at $t=0$ a tiny amount of noise remains
(which is why sampling ends with an explicit denoise at $t=0$ rather than just
reading off $\bz_0$), and at $t=T$ a tiny amount of signal remains, keeping
$\alpha_T>0$ so the division in the sampler is safe.

Measured values, with the signal-to-noise ratio $\text{SNR}=\alpha_t^2/\sigma_t^2$:

\begin{center}
\begin{tabular}{@{}rrrr@{}}
\toprule
$t$ & $\alpha_t$ & $\sigma_t$ & SNR \\
\midrule
$0$    & $0.999995$ & $0.003162$ & $1.0\times10^{5}$ \\
$100$  & $0.989995$ & $0.141101$ & $49.2$ \\
$500$  & $0.749999$ & $0.661439$ & $1.29$ \\
$900$  & $0.190024$ & $0.981779$ & $0.0375$ \\
$1000$ & $0.003163$ & $0.999995$ & $1.0\times10^{-5}$ \\
\bottomrule
\end{tabular}
\end{center}

\subsection{Scaling of the two modalities}

\begin{lstlisting}
def normalize(self, x, h):   return x, h * self.h_scale     # h_scale = 0.25
def unnormalize(self, x, h): return x, h / self.h_scale
\end{lstlisting}

\noindent
Shapes are untouched; only the magnitude of \code{h} changes. Coordinates have a
per-component standard deviation of about $1.4$\,\r{A} in QM9, whereas a one-hot
entry is exactly $0$ or $1$ --- and it is a \emph{discontinuous} target. Left
unscaled, the type channels dominate the gradient and the model under-fits
geometry. Multiplying by $0.25$ puts the two modalities on comparable footing.
Because the argmax used at decoding time is invariant to a positive rescaling,
this choice affects only the loss balance, never the decoded label.

\subsection{The training loss}

\begin{lstlisting}[caption={\code{EquivariantDiffusion.loss}}]
b = x.shape[0]
x, h = self.normalize(x, h)                     # [B,N,3], [B,N,K]
x    = remove_mean(x, node_mask)                # project data into X_0

t_int = torch.randint(0, self.timesteps + 1, (b,))      # [B]  integers in {0..T}
alpha = self.alphas[t_int].view(b, 1, 1)        # [B] -> [B,1,1]
sigma = self.sigmas[t_int].view(b, 1, 1)        # [B] -> [B,1,1]

eps_x = sample_zero_com_noise(x.shape, node_mask, x.device)   # [B,N,3] in X_0
eps_h = torch.randn_like(h) * node_mask.unsqueeze(-1)          # [B,N,K]

z_x = alpha * x + sigma * eps_x                 # [B,1,1]*[B,N,3] -> [B,N,3]
z_h = alpha * h + sigma * eps_h                 # [B,1,1]*[B,N,K] -> [B,N,K]

pred_x, pred_h = self.dynamics(t_int.float() / self.timesteps, z_x, z_h, node_mask)

err  = ((pred_x - eps_x)**2).sum() + ((pred_h - eps_h)**2).sum()   # -> scalar
n_atoms = node_mask.sum()                                          # -> scalar
dof  = 3 * (n_atoms - b) + self.num_types * n_atoms                # -> scalar
return err / dof
\end{lstlisting}

\noindent
Shape observations:
\begin{itemize}[leftmargin=1.4em]
\item \code{alphas} is indexed by a \shape{B} integer tensor, giving a \shape{B}
      result, reshaped to \shape{B,1,1} so it broadcasts over both the atom and
      feature axes. Each molecule in the batch gets its \emph{own} timestep ---
      this is what makes one pass over the data cover the whole of $[0,T]$.
\item \code{t\_int.float() / self.timesteps} converts integer timesteps to the $[0,1]$
      scalars the network expects, shape \shape{B}.
\item Both squared errors reduce to a scalar over \emph{all} axes at once. Padded
      slots contribute exactly $0$: \code{pred} is masked by the network and
      \code{eps} is masked at construction.
\end{itemize}

\paragraph{The normalising constant.} The loss is the eps-prediction objective
with weight $w(t)=1$,
\begin{equation}
  \mathcal{L} \;=\;
  \frac{\sum_b \bigl\lVert \hat{\bm}_x^{(b)} - \bm_x^{(b)} \bigr\rVert^2
      + \bigl\lVert \hat{\bm}_h^{(b)} - \bm_h^{(b)} \bigr\rVert^2}
       {\underbrace{3\bigl(\textstyle\sum_b n_b - B\bigr)}_{\text{coordinate d.o.f.}}
      + \underbrace{K\textstyle\sum_b n_b}_{\text{type d.o.f.}}}\,.
\end{equation}
The $-B$ is the geometry of Section~\ref{sec:com} showing up in the arithmetic:
each molecule contributes $3(n_b-1)$ coordinate degrees of freedom, not $3n_b$,
because the centre-of-mass constraint removes three. Dividing by the true
dimension makes the reported number a per-dimension mean, comparable across
batches with different molecule sizes --- and hence comparable between the
training and validation logs.

\paragraph{What is omitted.} This is EDM's $L_{\text{simple}}$ only. There is no
variational bound / negative log-likelihood term, no $L_0$ reconstruction term,
and no $t$-dependent reweighting. That is the standard practical choice --- the
paper reports that $w(t)=1$ trains better samples than the true ELBO weighting ---
but it means this code cannot report likelihoods.

\section{Training: \file{train.py}}

\subsection{The loop}

\begin{lstlisting}[caption={The inner training step}]
for batch in train_loader:
    x, h, mask = (batch[k].to(device) for k in ("x", "h", "mask"))  # [B,N,3] [B,N,K] [B,N]
    loss = model.loss(x, h, mask)                                   # scalar
    opt.zero_grad(set_to_none=True)
    loss.backward()
    grad_norm = torch.nn.utils.clip_grad_norm_(model.parameters(), args.clip_grad)
    opt.step()
    ema.update(model)
\end{lstlisting}

\noindent
Four details carry weight here.

\begin{description}[leftmargin=0pt,style=nextline]
\item[Gradient clipping is not optional.] The observed gradient norm on the first
      iteration is $\approx 530$, falling to $\approx 100$ within one epoch.
      Clipping at $1.5$ is what keeps the early updates from destroying the
      coordinate head. The norm is logged so this can be watched.
\item[AdamW with \code{amsgrad}, $\text{lr}=10^{-4}$, weight decay $10^{-12}$.]
      The decay is effectively off; it is present because EDM uses this exact
      configuration, and diffusion training is known to be sensitive to
      regularisation that fights the noise-prediction objective.
\item[EMA of the weights.] Diffusion sampling chains $T$ forward passes, so
      small weight noise compounds. \code{EMA} keeps a shadow copy updated as
      $\theta_{\text{ema}} \leftarrow 0.999\,\theta_{\text{ema}} + 0.001\,\theta$
      after every step, and it is the shadow that is sampled from. Its horizon
      is $\sim\!1000$ steps, under one epoch at $1562$ iterations/epoch.
\item[The shadow is a whole diffusion model.] \code{EMA} deep-copies the
      \code{EquivariantDiffusion} object, so \code{ema.shadow.sample(...)} works
      directly. The averaging loop iterates over \code{state\_dict()} values and
      skips non-floating-point entries, so integer buffers are copied rather
      than blended.
\end{description}

\subsection{Per-epoch evaluation}

Each epoch records train loss, validation loss and wall time; every
\code{--eval-every} epochs it additionally \emph{generates} molecules and scores
them. The size-to-mask conversion is the interesting shape step:

\begin{lstlisting}[caption={\code{sample\_and\_score}}]
sizes = torch.multinomial(size_hist, n_samples, replacement=True)      # [B]
mask  = (torch.arange(max_n)[None, :] < sizes[:, None]).float()        # [B,N]
x, one_hot, _ = model.sample(mask)                                     # [B,N,3], [B,N,K]
mols = to_molecules(x, one_hot, mask, symbols)      # dense -> list of (n_i,3) + symbols
return evaluate(mols, implicit_h), mols
\end{lstlisting}

\noindent
\code{to\_molecules} performs the inverse of the dataset's padding: it argmaxes
the one-hot back to indices \shape{B,N}, converts the mask to a boolean, and uses
it to slice each molecule down to its true length, yielding a ragged Python list
of $(n_i,3)$ arrays with matching element symbols. Checkpoints store both raw and
EMA weights, the config, and \code{size\_hist}, so \code{sample.py} can
reconstruct the model without being told the architecture.

\section{Sampling: \code{EquivariantDiffusion.sample}}
\label{sec:sampling}

\subsection{The reverse process}

Sampling runs the chain backwards from pure noise. The generative model is
factorised as $p(\bz_0) = p(\bz_T)\prod_{t=1}^{T} p(\bz_{t-1}\mid\bz_t)$, and
each reverse step is obtained by taking the true posterior
$q(\bz_s\mid\bz_t,\bz_0)$ for $s=t-1$ and substituting the network's estimate of
$\bz_0$. Define
\begin{equation}
  \alpha_{t|s} = \frac{\alpha_t}{\alpha_s},
  \qquad
  \sigma^2_{t|s} = \sigma_t^2 - \alpha_{t|s}^2\,\sigma_s^2 .
\end{equation}
The Gaussian posterior is
\begin{equation}
  q(\bz_s\mid \bz_t,\bz_0)
  = \mathcal{N}\!\Bigl(
      \tfrac{\alpha_{t|s}\sigma_s^2}{\sigma_t^2}\,\bz_t
      + \tfrac{\alpha_s \sigma^2_{t|s}}{\sigma_t^2}\,\bz_0,
      \;\;
      \tfrac{\sigma^2_{t|s}\,\sigma_s^2}{\sigma_t^2}\,\mathbf{I}
    \Bigr).
\end{equation}
Substituting $\bz_0=(\bz_t-\sigma_t\bm)/\alpha_t$ and collecting the $\bz_t$
terms,
\begin{equation}
  \alpha_{t|s}\sigma_s^2 + \frac{\sigma^2_{t|s}}{\alpha_{t|s}}
  = \frac{\alpha_{t|s}^2\sigma_s^2 + \sigma_t^2 - \alpha_{t|s}^2\sigma_s^2}{\alpha_{t|s}}
  = \frac{\sigma_t^2}{\alpha_{t|s}},
\end{equation}
so the mean collapses to a two-term expression in $\bz_t$ and the predicted
noise, which is exactly what the code computes:
\begin{equation}
  \boxed{\;
  \boldsymbol{\mu}_{t\to s} \;=\; \frac{\bz_t}{\alpha_{t|s}}
    \;-\; \frac{\sigma^2_{t|s}}{\alpha_{t|s}\,\sigma_t}\,\hat{\bm}\,,
  \qquad
  \text{std} \;=\; \frac{\sigma_{t|s}\,\sigma_s}{\sigma_t}\;}
\end{equation}

\begin{lstlisting}[caption={\code{\_p\_step} --- one ancestral step $\bz_t\to\bz_{t-1}$}]
s_int = t_int - 1
a_t, s_t = self.alphas[t_int], self.sigmas[t_int]     # 0-dim tensors (scalars)
a_s, s_s = self.alphas[s_int], self.sigmas[s_int]

a_ts   = a_t / a_s
var_ts = (s_t**2 - a_ts**2 * s_s**2).clamp(min=1e-12)
std    = (var_ts.sqrt() * s_s / s_t).view(1, 1, 1)    # -> [1,1,1], broadcasts

t = torch.full((b,), t_int / self.timesteps)                  # [B]
pred_x, pred_h = self.dynamics(t, z_x, z_h, node_mask)   # [B,N,3], [B,N,K]

coef = (var_ts / (a_ts * s_t)).view(1, 1, 1)          # [1,1,1]
mu_x = z_x / a_ts - coef * pred_x                     # [B,N,3]
mu_h = z_h / a_ts - coef * pred_h                     # [B,N,K]

z_x = mu_x + std * sample_zero_com_noise(z_x.shape, node_mask, z_x.device)
z_h = mu_h + std * torch.randn_like(z_h) * node_mask.unsqueeze(-1)
return remove_mean(z_x, node_mask), z_h * node_mask.unsqueeze(-1)
\end{lstlisting}

\noindent
Unlike training, sampling uses \emph{one} timestep for the whole batch, so
$\alpha,\sigma$ are $0$-dimensional tensors reshaped to \shape{1,1,1} rather
than \shape{B,1,1}. The coordinate noise is again drawn in $\mathcal{X}_0$, and
\code{remove\_mean} is reapplied at the end of every step so floating-point
drift cannot accumulate over $1000$ iterations.

\subsection{The chain and the final decode}

\begin{lstlisting}[caption={\code{sample} --- the full reverse trajectory}]
z_x = sample_zero_com_noise((b, n, 3), node_mask, device)          # [B,N,3]
z_h = torch.randn(b, n, self.num_types, device=device) * nm        # [B,N,K]

for t_int in range(self.timesteps, 0, -1):                                 # T .. 1
    z_x, z_h = self._p_step(z_x, z_h, t_int, node_mask)

t = torch.zeros(b, device=device)                                  # [B]
pred_x, pred_h = self.dynamics(t, z_x, z_h, node_mask)
a0, s0 = self.alphas[0], self.sigmas[0]
x = remove_mean((z_x - s0 * pred_x) / a0, node_mask)               # [B,N,3]
h = (z_h - s0 * pred_h) / a0                                       # [B,N,K]

x, h    = self.unnormalize(x, h)                                   # h /= 0.25
types   = h.argmax(-1)                                             # [B,N,K] -> [B,N]
one_hot = F.one_hot(types, self.num_types).float() * nm             # [B,N] -> [B,N,K]
return x * nm, one_hot, frames
\end{lstlisting}

\noindent
The loop runs $T$ times and each iteration is one full network forward pass ---
sampling costs $1000\times$ a training forward pass, which is why generating a
few hundred molecules takes longer than an epoch of training on them.

Two steps close the process:
\begin{enumerate}[leftmargin=1.6em]
\item \textbf{A final denoise at $t=0$.} Because $\sigma_0=0.0032\neq0$, $\bz_0$
      is still very slightly noisy; $\bx = (\bz_x - \sigma_0\hat{\bm}_x)/\alpha_0$
      removes that last increment.
\item \textbf{Discretisation of the atom types.} \code{argmax} over the last axis
      collapses \shape{B,N,K} to \shape{B,N} integer labels, which are re-expanded
      to a clean one-hot. This is the point at which the continuous relaxation
      becomes a discrete chemical decision --- and where information is
      irrecoverably thrown away, since a near-tie between carbon and nitrogen is
      resolved silently.
\end{enumerate}

\subsection{A closed-form check of the sampler}
\label{sec:calib}

The reverse-process algebra above can be validated without any trained network.
If the data were $\bx\sim\mathcal{N}(0,\mathbf{I})$, then since
$\bz_t\sim\mathcal{N}(0,(\alpha_t^2+\sigma_t^2)\mathbf{I})=\mathcal{N}(0,\mathbf{I})$,
the optimal denoiser is available in closed form:
\begin{equation}
  \E[\bm \mid \bz_t] \;=\; \sigma_t\,\bz_t .
\end{equation}
Substituting this exact denoiser for the network must return samples with unit
variance --- any error in $\alpha_{t|s}$, $\sigma^2_{t|s}$ or the standard
deviation shows up as a drift in scale. The test in \code{test\_smoke.py} does
precisely this with $n=6$ atoms and recovers a standard deviation of
$\mathbf{0.911}$ against the analytic prediction for the projected subspace,
\begin{equation}
  \sqrt{\tfrac{n-1}{n}} \;=\; \sqrt{5/6} \;=\; 0.913 .
\end{equation}
The agreement, including the subspace correction, is the evidence that the
schedule and the ancestral step are implemented correctly.

\section{Evaluation: \file{qm9\_diffusion/stability.py}}
\label{sec:eval}

The model outputs coordinates, not bonds, so the metric has to infer chemistry
from geometry. For each unordered pair, the interatomic distance is compared
against tabulated single/double/triple bond lengths for that element pair, with
tolerances of $(10,5,3)$\,pm; the resulting bond orders are summed per atom to
give an inferred valence.

\begin{center}
\begin{tabular}{@{}lll@{}}
\toprule
\textbf{Metric} & \textbf{Definition} & \textbf{Shape reduction} \\
\midrule
atom stability     & inferred valence $=$ expected valence & $(n,3)\to n$ booleans $\to$ scalar \\
molecule stability & \emph{all} atoms in the molecule stable & $n$ booleans $\to$ 1 boolean \\
\bottomrule
\end{tabular}
\end{center}

\noindent
Expected valences are H:1, C:4, N:3, O:2, F:1. Molecule stability is the harsher
metric by construction --- one wrong atom in $18$ fails the whole molecule --- so it
lags atom stability substantially in an undertrained model.

\paragraph{Validating the metric itself.} Run on \emph{real} QM9 geometries the
implementation gives $0.994$ atom / $0.951$ molecule stability, against the
published $0.99$/$0.95$. Reproducing the reference numbers on reference data is
what establishes that the tolerances and bond tables are right, so that a poor
score on generated molecules can be attributed to the model.

\paragraph{The heavy-atom caveat.} With \code{--remove-h}, exact valence is
unsatisfiable: a carbon that came from methane has no heavy-atom neighbours at
all. In that mode the check is relaxed to \emph{valence $\le$ allowed}, treating
the deficit as implicit hydrogens. On real data this reads $0.996$/$0.978$. It
is a much weaker test, and no-H numbers must not be compared with with-H ones.

\section{Masking: the invariant to preserve}
\label{sec:masking}

Because the implementation is dense, correctness rests on padded slots never
influencing anything. The places that enforce it:

\begin{center}\small
\begin{tabular}{@{}llp{5.6cm}@{}}
\toprule
\textbf{Location} & \textbf{Mechanism} & \textbf{What breaks without it} \\
\midrule
dataset \code{\_\_getitem\_\_} & \code{* mask[:, None]} & padded rows decode as spurious hydrogens \\
\code{remove\_mean}    & masked mean, divide by $n_b$ & centre of mass computed over $N$, not $n_b$ \\
\code{edge\_mask}      & outer product of masks, minus $I$ & phantom atoms send messages \\
\code{n\_edges}        & \code{edge\_mask.sum(2)}   & aggregation normalised by $N-1$ \\
every layer output     & \code{* node\_mask}        & padding accumulates state across layers \\
loss                   & \code{eps} masked at creation & padded slots contribute error \\
\code{dof}             & \code{node\_mask.sum()}    & per-dimension mean depends on padding \\
\bottomrule
\end{tabular}
\end{center}

\noindent
\code{test\_masking\_and\_com} asserts the output invariant directly: for a batch
containing molecules of sizes $9,7,5,3$ padded to $N=9$, both output tensors are
$0$ in every padded slot, and the masked centre of mass is zero to
$1.5\times10^{-11}$.

\section{Computational cost}
\label{sec:cost}

\subsection{Where the memory goes}

The dense formulation materialises $B\!\times\!N^2$ edges, each carrying $H$
channels. For $B=64$, $N=29$, $H=192$ in \code{float32}:
\begin{center}
\begin{tabular}{@{}lrr@{}}
\toprule
\textbf{Tensor} & \textbf{Elements} & \textbf{Memory} \\
\midrule
\code{edge\_mlp} input \shape{B,N,N,2H{+}1} & $20.7$\,M & $82.9$\,MB \\
message \code{m} \shape{B,N,N,H}            & $10.3$\,M & $41.3$\,MB \\
node tensor \shape{B,N,H}                   & $0.36$\,M & $1.4$\,MB \\
\bottomrule
\end{tabular}
\end{center}
Several such activations are retained per layer for the backward pass, so
training memory scales as $\mathcal{O}(L\cdot B\cdot N^2\cdot H)$ --- the edge
tensor, not the parameters ($1.6$\,M, about $6$\,MB), is what fills the device.
Reduce \code{--batch-size} or \code{--hidden} first when memory is tight.

For QM9 the dense choice costs almost nothing: a fully connected $29$-atom
molecule has $812$ directed edges, and a sparse implementation of the same model
would build essentially the same number. Dense wins on simplicity here; it would
not for larger molecules.

\subsection{Measured throughput}

On Apple Silicon (MPS), batch size $64$, one epoch $=100$k molecules:

\begin{center}
\begin{tabular}{@{}lrrr@{}}
\toprule
\textbf{Configuration} & \textbf{Params} & \textbf{ms/iter} & \textbf{min/epoch} \\
\midrule
$H{=}256$, $L{=}9$ (EDM paper) & $4.22$\,M & $613$ & $16.0$ \\
$H{=}192$, $L{=}6$ (default)   & $1.60$\,M & $294$ & $7.7$ \\
$H{=}128$, $L{=}4$             & $0.48$\,M & $123$ & $3.2$ \\
$H{=}128$, $L{=}4$, \code{--remove-h} & $0.48$\,M & $21$ & $0.5$ \\
\bottomrule
\end{tabular}
\end{center}

\noindent
The last two rows are the same model on the same number of molecules; only $N$
changed, from $29$ to $9$. The measured speedup is $123/21 \approx 5.9\times$,
against the $\left(29/9\right)^2 \approx 10.4\times$ that pure $N^2$ scaling of
the edge axis would predict --- the shortfall is per-iteration overhead
(kernel launches, the optimiser step, host-device synchronisation) that does not
shrink with $N$. It is nevertheless the reason the heavy-atom mode is the
practical setting for a laptop.

\section{Summary: one forward pass end to end}

Collecting every shape change for a single training step with the default
configuration $B{=}64$, $N{=}29$, $K{=}5$, $H{=}192$, $L{=}6$:

\begin{center}
\begin{tabular}{@{}rlll@{}}
\toprule
\textbf{\#} & \textbf{Operation} & \textbf{Shape} & \textbf{Where} \\
\midrule
1  & batch from \code{DataLoader}      & \shape{64,29,3}, \shape{64,29,5}, \shape{64,29} & \code{data.py} \\
2  & scale types by $0.25$             & unchanged & \code{normalize} \\
3  & project coordinates into $\mathcal{X}_0$ & unchanged & \code{remove\_mean} \\
4  & draw one timestep per molecule    & \shape{64} & \code{loss} \\
5  & gather $\alpha_t,\sigma_t$        & \shape{64} $\to$ \shape{64,1,1} & \code{loss} \\
6  & add noise                         & \shape{64,29,3}, \shape{64,29,5} & \code{loss} \\
\addlinespace
7  & embed types $+$ time              & \shape{64,29,6} $\to$ \shape{64,29,192} & \code{embed} \\
8  & build the edge mask               & \shape{64,29,1} $\to$ \shape{64,29,29,1} & \code{forward} \\
9  & $\times L{=}6$: lift to edges     & \shape{64,29,29,385} & \code{edge\_mlp} \\
10 & \phantom{$\times L$:} message     & \shape{64,29,29,192} & \code{edge\_mlp} \\
11 & \phantom{$\times L$:} coordinate scalar & \shape{64,29,29,1} & \code{coord\_mlp} \\
12 & \phantom{$\times L$:} contract to nodes & \shape{64,29,192}, \shape{64,29,3} & sum over axis 2 \\
\addlinespace
13 & coordinate displacement, re-projected & \shape{64,29,3} & \code{forward} \\
14 & decode type noise                 & \shape{64,29,192} $\to$ \shape{64,29,5} & \code{decode} \\
15 & squared error / d.o.f.            & $\to$ scalar & \code{loss} \\
\bottomrule
\end{tabular}
\end{center}

\noindent
At sampling time steps 7--14 are repeated $T=1000$ times with the shapes
unchanged, followed by an \code{argmax} that collapses \shape{B,N,K} to
\shape{B,N} and a final one-hot re-expansion.

\section*{References}
\begin{enumerate}[leftmargin=1.6em]
\item E.~Hoogeboom, V.~G.~Satorras, C.~Vignac, M.~Welling.
      \emph{Equivariant Diffusion for Molecule Generation in 3D}. ICML 2022.
      \url{https://arxiv.org/abs/2203.17003}
\item V.~G.~Satorras, E.~Hoogeboom, M.~Welling.
      \emph{E(n) Equivariant Graph Neural Networks}. ICML 2021.
      \url{https://arxiv.org/abs/2102.09844}
\item J.~Ho, A.~Jain, P.~Abbeel.
      \emph{Denoising Diffusion Probabilistic Models}. NeurIPS 2020.
      \url{https://arxiv.org/abs/2006.11239}
\item R.~Ramakrishnan, P.~O.~Dral, M.~Rupp, O.~A.~von Lilienfeld.
      \emph{Quantum chemistry structures and properties of 134 kilo molecules}.
      Scientific Data 1, 140022 (2014).
\end{enumerate}

\end{document}
